\section{Evaluation}
\label{sec:evaluation}

This section evaluates CPSForge as a closed-loop LLM red/blue team framework across three Factory~I/O scenes on a real Siemens S7-1200 PLC. All reported numbers are collected from live PLC runs; no results are simulated or mocked.

\subsection{Research Questions}

\begin{description}[leftmargin=1.5em, style=nextline]
    \item[RQ1 (Context Ablation):] How does operational context (minimal vs.\ partial vs.\ full) affect the capability of a local LLM to launch successful CPS attacks?
    \item[RQ2 (Fine-Tuning):] How does task-specific QLoRA fine-tuning change attack capability, stealth, validity, and cross-scene transfer on Qwen2.5-3B-Instruct?
    \item[RQ3 (Attack Mode):] Can an online MITM LLM attacker exploit timing and phase transitions more effectively than a static one-shot LLM or a random baseline?
    \item[RQ4 (Defenses):] Which runtime defenses remain effective against context-aware, timing-aware LLM attackers? Does an LLM defender outperform rule-based defenses?
    \item[RQ5 (Cross-Model Robustness):] Are the observed attack and defense trends robust across three local models from two families, or are they strongly model-specific?
    \item[RQ6 (Frontier Model):] Does a frontier API model (GPT-4o-mini) exhibit qualitatively different attack capability, self-deterrence behavior, or defense evasion compared to local 3B-class models?
\end{description}

\subsection{Evaluation Environment}

Experiments are conducted on a real Siemens S7-1200 PLC connected to Factory~I/O~\cite{factoryio} through TIA Portal~v17~\cite{tiaportal}. The evaluation covers three scenes: Level Control (continuous tank process, 15~tags), Sorting by Height (discrete 2-way height sorting, 30~tags), and Sorting by Weight (discrete 3-way weight sorting, 35~tags). The middleware polls PLC tags at 500\,ms intervals via python-snap7~\cite{snap7}; the LLM attacker is called every three polling cycles. We note that the S7-1200 PLC scan cycle is typically 1--10\,ms; the 500\,ms polling interval is a middleware-level parameter that determines observation granularity, not PLC control speed. Attack effects that occur and are corrected within a single polling interval may not be observable.

All LLM inference runs in-process via HuggingFace \texttt{transformers} with 4-bit NF4 quantization through \texttt{bitsandbytes}. No external inference server is used. The primary model is Qwen2.5-3B-Instruct (${\sim}$6\,GB), studied in base form. Two additional cross-model comparison models are used for RQ5: Qwen3-1.7B (${\sim}$4\,GB) and SmolLM3-3B (${\sim}$6\,GB). Experiments run on a Dell Precision 5820 with an NVIDIA RTX A4000 (16\,GB VRAM) under Windows~11 Enterprise.

The completed experiment matrix comprises 333~runs: 27~(RQ1) + 36~(RQ2) + 27~(RQ3) + 126~(RQ4, base + finetuned) + 72~(RQ5) + 45~(RQ6). Each run executes 60~steps at 500\,ms polling (30\,s of PLC polling time). Including LLM inference---13--33\,s per call at up to 20~calls per run for local models, 2--5\,s for the API model---total wall-clock time is approximately 1--10~minutes per run. Each run produces a unified per-step log (Parquet with 33+ fields), attack records, shield decisions, PLC event captures, and summary metrics (all JSON).

\subsection{Attacker Configurations}

We compare three attacker modes (RQ3), each producing the same unified per-step log format:

\begin{itemize}[leftmargin=1.5em]
    \item \textbf{Random baseline:} bounded random perturbations over the configured attack surface, providing a lower bound on attacker capability. Always attacks (no wait decisions).
    \item \textbf{Static LLM:} a one-shot batch attacker that receives full context once and generates attack proposals without observing process feedback. Tests LLM attack capability without online adaptation.
    \item \textbf{Online MITM:} the core contribution (C1). An iterative observe--decide--act loop where the LLM observes live plant state every three control cycles, infers the current process phase, builds a context-tier-specific prompt, and decides to \emph{attack} or \emph{wait}. All writes pass through the full defense chain before reaching the PLC.
\end{itemize}

\subsection{Defense Configurations}

We evaluate seven defense variants for RQ4:

\begin{itemize}[leftmargin=1.5em]
    \item \textbf{None:} no defense (measures raw attacker capability).
    \item \textbf{Baseline:} threshold and invariant detectors.
    \item \textbf{Phase-aware shield (C4):} blocks writes inconsistent with the inferred process phase.
    \item \textbf{Intent-consistency checker (C4):} blocks writes opposing the controller's current trend.
    \item \textbf{Combined:} phase-aware shield + intent-consistency checker together.
    \item \textbf{LLM defender (novel):} an LLM-based anomaly scorer using the same model and context pipeline as the attacker. Evaluates each proposed write and decides allow/block with a suspicion score.
    \item \textbf{LLM combined:} phase-aware + intent-consistency + LLM defender together.
\end{itemize}

\subsection{Metrics}

\textbf{Attacker metrics.} Attack Success Rate (ASR) $= N_{\text{succ}}/N_{\text{exec}}$, computed post-hoc: for write at step $N$, compare tag value at step $N{-}1$ vs.\ $N{+}2$; success if the tag moved toward the injected value. Valid Action Rate (VAR) $= N_{\text{schema\_valid}}/N_{\text{submitted}}$, and mean LLM inference latency (ms).

\textbf{Defense metrics.} Prevention rate $= N_{\text{blocked}}/N_{\text{proposals}}$, measuring the fraction of attack proposals blocked before PLC write. Includes safety-shield blocks; rows where ``None'' defense shows nonzero prevention reflect mandatory safety-shield blocking, not RQ4 defense activation. Per-stage block counts (Ph., Int., LLM columns) identify the specific defense stage that blocked each proposal.

\textbf{Statistical note.} All tables report 3-repeat cells. For rare-event metrics (ASR), point estimates carry wide uncertainty (95\% CI width $\approx \pm 30\,\text{pp}$ for $n \leq 10$ writes). Findings are supported where multiple independent cells show the same pattern.

\subsection{Results}

\textbf{Safety-shield caveat.} All reported ASR values are \emph{lower bounds} under the experimental safety envelope: the mandatory safety shield blocks out-of-range and invariant-violating writes before they reach the PLC. An unconstrained adversary with the same S7 access could execute writes that the shield rejects, potentially achieving higher ASR. This constraint applies uniformly to all cells and does not affect within-table comparisons, but absolute ASR values should not be interpreted as the ceiling of LLM attacker capability.

\subsubsection{RQ1: Context Ablation}
\label{sec:eval:rq1}

Table~\ref{tab:context_ablation} shows the effect of operational context on attack capability using Qwen2.5-3B-Instruct (base) as the online MITM attacker with no defense active.

% Table 3 --- RQ1 Context Ablation
% Generated by: python -m cpsforge.analysis.generate_tables
% RQ-specific filtering: limited to 3 runs per (scene, context) cell
\begin{table}[t]
\centering
\caption{RQ1: Context ablation for online MITM attacker (base Qwen2.5-3B; no RQ4 defense active; safety shield always active).}
\label{tab:context_ablation}
\small
\setlength{\tabcolsep}{3pt}
\begin{tabular}{llrrrrrr}
\toprule
\textbf{Scene} & \textbf{Context} & \textbf{Runs} & \textbf{Prop.} & \textbf{Writes} & \textbf{Succ.} & \textbf{ASR\%} & \textbf{Lat.(ms)} \\
\midrule
Level Control & Minimal & 3 & 23 & 2 & 0 & 0.0 & 13086 \\
 & Partial & 3 & 20 & 13 & 0 & 0.0 & 18429 \\
 & Full & 3 & 5 & 0 & 0 & 0.0 & 7863 \\
\midrule
Sort.~Weight & Minimal & 3 & 13 & 13 & 7 & 53.8 & 17929 \\
 & Partial & 3 & 12 & 5 & 2 & 40.0 & 19789 \\
 & Full & 3 & 0 & 0 & 0 & 0.0 & 11180 \\
\midrule
Sort.~Height & Minimal & 3 & 30 & 30 & 1 & 3.3 & 16906 \\
 & Partial & 3 & 30 & 25 & 1 & 4.0 & 22944 \\
 & Full & 3 & 30 & 30 & 4 & 13.3 & 22263 \\
\bottomrule
\end{tabular}
\vspace{2pt}
\noindent\footnotesize{3 repeats per configuration; ASR computed post-hoc.}
\end{table}

\textbf{Finding 1: Full context causes self-deterrence on Sorting by Weight and near-complete suppression on Level Control.} Full context produces \emph{zero} attack proposals on Sorting by Weight---the model decides ``wait'' on every LLM call across all 3~repeats. On Level Control, full context sharply reduces proposals from 23~(minimal) to just 5, all of which the safety shield blocks (Writes~=~0). The rich operational context (phase inference, safety rules, control objectives) makes the 3B model overly cautious: on Sort.~Weight it fully self-deterres; on Level Control it mostly self-deterres and the few proposals it does generate use out-of-range values that the shield rejects.

\textbf{Finding 2: Self-deterrence is scene-dependent.} On Sorting by Weight, partial context achieves 40.0\% ASR vs.\ 53.8\% at minimal---both are effective, while full context drops to zero proposals (complete self-deterrence). On Level Control, all three context levels produce zero successful attacks; full context reduces proposals to 5 (all shield-blocked) while minimal and partial generate 23 and 20 proposals respectively. On Sorting by Height, full context is the \emph{most} effective tier (13.3\% ASR vs.\ 3.3\% and 4.0\% at minimal and partial), suggesting that the scene's discrete state-machine logic benefits from richer phase and trend context.

\textbf{Finding 3: Context changes attack strategy, not just success rate.} With minimal context, the model uses brute-force actuator overrides. With partial context, it shifts toward control-plane setpoint manipulation---a more targeted strategy. Full context, when it does produce attacks on Sorting by Height, generates exactly 30~proposals all executed (13.3\% ASR), targeting a narrow set of tags with deliberate actuator overrides. More context makes the model more selective but also more conservative; the direction of the net effect depends on the scene's control paradigm (continuous PID vs.\ discrete state machine).

\textbf{Note on Level Control results.} Level Control shows 0\% ASR at all three context tiers: minimal (0/2 writes succeed), partial (0/13 writes succeed), and full (5 proposals, all shield-blocked, 0 writes executed). The partial-context model generates many proposals that pass the safety shield but fail to shift the PLC's PID-controlled tag, suggesting that the PID correction loop counteracts the attacker's writes faster than the 500\,ms polling window can observe. The full-context model proposes values outside the configured safety range, so its proposals never reach the PLC.

\textbf{Implication for defense:} The self-deterrence effect suggests that \emph{exposing operational context may itself be a defense mechanism}---a form of deterrence through transparency that reduces LLM red team effectiveness without any active blocking.

\subsubsection{RQ2: Fine-Tuning Impact}
\label{sec:eval:rq2}

RQ2 measures how task-specific QLoRA fine-tuning changes \emph{both} red team (attack) and blue team (defense) capability on the same model. The training dataset (2{,}057 examples) includes: (i)~83 post-hoc verified successful attack proposals; (ii)~608 defense ``block'' examples from known attack proposals; and (iii)~1{,}366 defense ``allow'' examples from normal operation states. Fine-tuning uses QLoRA (rank=16, $\alpha$=32) on Qwen2.5-3B-Instruct in 4-bit NF4 quantization for 3~epochs.

\textbf{Training procedure.} The 2{,}057 examples are used in a single training split without a held-out validation set, following the approach of Hackphyr~\cite{rigaki2024hackphyr} (1{,}641 examples, no validation split). We monitor training loss, which decreases monotonically from 1.82 to 0.41 over 3~epochs. The absence of a validation split means we cannot formally rule out overfitting; however, the adapter's total parameter count (29.9M, 0.96\% of the 3.1B base) and LoRA's implicit regularization (low-rank weight updates) inherently limit overfitting capacity. We train with a single seed; multi-seed training would strengthen confidence in the results but was not feasible given hardware constraints. The dataset is small relative to comparable cybersecurity fine-tuning work: HackMentor uses 14{,}000 examples and CyberLLMInstruct uses 54{,}928, though our task is narrower (structured JSON attack/defense decisions for three specific CPS scenes rather than general cybersecurity Q\&A). The 83 verified attack examples are the primary limitation---this is the number of physically successful writes observed across 189~base-model runs, and cannot be synthetically expanded without risking distributional artifacts.

Table~\ref{tab:finetune_comparison} shows the results.

% Table RQ2 --- Fine-Tuning Comparison
% Generated by: python -m cpsforge.analysis.generate_tables
% RQ-specific filtering: limited to 3 runs per (scene, context, model) cell
\begin{table}[t]
\centering
\caption{RQ2: Fine-tuning effect on attack capability (online MITM, no defense). Base = Qwen2.5-3B-Instruct; FT = same model with QLoRA v2 adapter. VAR = schema-valid outputs / LLM calls.}
\label{tab:finetune_comparison}
\small
\setlength{\tabcolsep}{3pt}
\resizebox{\columnwidth}{!}{%
\begin{tabular}{lllrrrrrrr}
\toprule
\textbf{Scene} & \textbf{Ctx} & \textbf{Model} & \textbf{Runs} & \textbf{LLM calls} & \textbf{VAR\%} & \textbf{Prop.} & \textbf{Writes} & \textbf{Succ.} & \textbf{ASR\%} \\
\midrule
Level Control & Minimal & Base & 3 & 27 & 100.0 & 23 & 2 & 0 & 0.0 \\
 & Minimal & FT & 3 & 30 & 100.0 & 30 & 30 & 11 & 36.7 \\
 & Full & Base & 3 & 15 & 100.0 & 5 & 0 & 0 & 0.0 \\
 & Full & FT & 3 & 34 & 0.0 & 0 & 0 & 0 & 0.0 \\
\midrule
Sort.~Weight & Minimal & Base & 3 & 14 & 92.9 & 13 & 13 & 7 & 53.8 \\
 & Minimal & FT & 3 & 30 & 100.0 & 30 & 16 & 0 & 0.0 \\
 & Full & Base & 3 & 5 & 40.0 & 0 & 0 & 0 & 0.0 \\
 & Full & FT & 3 & 60 & 0.0 & 0 & 0 & 0 & 0.0 \\
\midrule
Sort.~Height & Minimal & Base & 3 & 30 & 100.0 & 30 & 30 & 1 & 3.3 \\
 & Minimal & FT & 3 & 30 & 100.0 & 30 & 30 & 3 & 10.0 \\
 & Full & Base & 3 & 30 & 100.0 & 30 & 30 & 4 & 13.3 \\
 & Full & FT & 3 & 45 & 0.0 & 0 & 0 & 0 & 0.0 \\
\bottomrule
\end{tabular}%
}
\vspace{2pt}
\noindent\footnotesize{3 repeats per configuration; ASR computed post-hoc.}
\end{table}

\textbf{Finding RQ2-1: Fine-tuning substantially improves attack effectiveness on Level Control minimal.} The fine-tuned model generates 30~proposals, all of which execute (30 writes, 11 successes, 36.7\% ASR), compared to the base model's 23 proposals but only 2~executed writes and 0\% ASR. Task-specific training enables the model to propose values that pass the safety shield---where the base model frequently proposes out-of-range values that are blocked---and succeed in shifting the target tag.

\textbf{Finding RQ2-2: Fine-tuning does not improve all scenes uniformly.} On Sorting by Weight minimal, fine-tuning \emph{hurts} attack effectiveness: the fine-tuned model generates more proposals (30 vs.\ 13) but achieves 0\% ASR compared to 53.8\% for the base. The adapter has learned to propose more actions but cannot replicate the base model's scene-specific targeting on this 35-tag complex scene.

\textbf{Finding RQ2-3: Fine-tuned model produces zero proposals under full context due to schema regression.} On all three scenes under full context, the fine-tuned model makes LLM calls (34--60 per scene) but achieves VAR~=~0.0\%---every output fails schema validation. This is \emph{schema regression}, not strategic self-deterrence: the base model under full context achieves VAR~=~100\% on Level Control and Sort.~Height (issuing valid ``wait'' decisions) and VAR~=~40\% on Sort.~Weight (already struggling with the 35-tag prompt length), whereas the fine-tuned model achieves VAR~=~0\% on all three---it cannot produce parseable JSON under the longer full-context prompts. We attribute this to a \emph{prompt-length distributional mismatch}: the fine-tuning examples are derived from partial- and minimal-context runs (which produced successful attacks), so the adapter's training distribution is biased toward shorter prompts (${\sim}$400--800~tokens). Full-context prompts (${\sim}$1{,}500--2{,}000~tokens) fall outside this distribution, causing the adapter to generate malformed output. This is a known failure mode of LoRA adapters trained on narrow prompt distributions~\cite{hu2025evodefender}; mitigations include augmenting training data with full-context examples or increasing LoRA rank. Under minimal context, where prompt length aligns with the training distribution, fine-tuning raises Sort.~Height ASR from 3.3\% to 10.0\%, confirming the adapter is effective within its training distribution.

\textbf{Mitigation note.} The schema regression under full context is addressable: augmenting the v2 training dataset with ${\sim}$200 full-context examples (both attack and defense) would align the adapter's prompt-length distribution with all three context tiers. We did not implement this mitigation within the current evaluation timeline but note it as a straightforward extension for future work.

\textbf{Finding RQ2-4: The dual-role hypothesis holds for defense (see RQ4 finetuned).} Table~\ref{tab:defense_finetuned} (Section~\ref{sec:eval:rq4}) shows that the fine-tuned LLM defender achieves 100\% prevention on Sorting by Height (30/30 proposals blocked by the LLM component), identical to the base model's LLM defender. Training on both attack and defense examples preserves defense capability while modifying attack behavior---the fine-tuned model that achieves 36.7\% ASR in attack mode still blocks 100\% of proposals in defense mode on the same scene type.

\subsubsection{RQ3: Online vs.\ Static Attack Mode}
\label{sec:eval:rq3}

Table~\ref{tab:attacker_comparison} compares the three attacker modes across all three scenes with no defense active.

% Table 4 --- RQ3 Attacker Comparison
% Generated by: python -m cpsforge.analysis.generate_tables
% RQ-specific filtering: limited to 3 runs per (scene, attacker) cell
\begin{table}[t]
\centering
\caption{RQ3: Attacker comparison (base Qwen2.5-3B, no defense). Random = bounded random perturbations; Static LLM = one-shot batch; Online MITM = core contribution.}
\label{tab:attacker_comparison}
\small
\setlength{\tabcolsep}{3pt}
\begin{tabular}{llrrrrr}
\toprule
\textbf{Scene} & \textbf{Attacker} & \textbf{Runs} & \textbf{Prop.} & \textbf{Writes} & \textbf{Succ.} & \textbf{ASR\%} \\
\midrule
Level Control & Random & 3 & 10 & 10 & 1 & 10.0 \\
 & Static LLM & 3 & 0 & 0 & 0 & 0.0 \\
 & Online MITM & 3 & 14 & 0 & 0 & 0.0 \\
\midrule
Sort.~Weight & Random & 3 & 0 & 0 & 0 & 0.0 \\
 & Static LLM & 3 & 0 & 0 & 0 & 0.0 \\
 & Online MITM & 3 & 13 & 13 & 7 & 53.8 \\
\midrule
Sort.~Height & Random & 3 & 30 & 30 & 4 & 13.3 \\
 & Static LLM & 3 & 0 & 0 & 0 & 0.0 \\
 & Online MITM & 3 & 30 & 30 & 1 & 3.3 \\
\bottomrule
\end{tabular}
\vspace{2pt}
\noindent\footnotesize{3 repeats per configuration; ASR computed post-hoc.}
\end{table}

\textbf{Finding 4: Static LLM generates zero valid proposals under the batch design.} The static one-shot LLM produces no schema-valid attack proposals across all scenes. The batch attacker requests a full JSON array of up to 10 attack objects in a single call with a 350-token output budget---insufficient for the full array. Even after truncation repair (salvaging any complete objects from a partial response), zero schema-valid proposals survive. This reflects two concurrent constraints: (i)~the 350-token limit is too small for a multi-object JSON array, and (ii)~without iterative plant feedback, the model produces low-confidence generic outputs that fail schema validation. \textbf{Important caveat:} the 0\% result reflects our specific experimental design (350-token budget, batch JSON array format) rather than a fundamental limitation of one-shot LLM attack capability. AttackLLM~\cite{ahmed2025attackllm} uses GPT-4 with much larger context and output windows plus human-in-the-loop validation; such a design would likely overcome the token-budget constraint. We report the static result as evidence that \emph{our} batch design is ineffective, not that all one-shot designs are inherently incapable. The random baseline, which operates without any LLM call, achieves 10.0\% on Level Control and 13.3\% on Sorting by Height, establishing that the 0\% static LLM result is not the floor for all non-online approaches.

\textbf{Finding 5: Online MITM is the only effective attacker on complex scenes.} On Sorting by Weight, the online MITM achieves 53.8\% ASR while both random and static attackers generate zero proposals. The LLM's ability to reason about the 35-tag scene's control logic and select appropriate targets is essential---random perturbations cannot identify valid attack surfaces on complex scenes.

\textbf{Finding 6: Random baseline is competitive on discrete scenes.} On Sorting by Height, the random baseline achieves 13.3\% ASR compared to 3.3\% for the online MITM (3 runs each). Discrete sorting scenes have binary actuators where random single-tag overrides can accidentally align with vulnerable states, while the LLM's cautious wait/attack reasoning sometimes misses opportunistic windows.

\textbf{Finding 7: Level Control is resistant to all attacker types.} On Level Control, random achieves 10.0\% ASR with 10 writes, and static LLM generates 0 proposals. The online MITM generates 14 proposals but all are blocked by the safety shield (Writes~=~0)---the LLM proposes out-of-range values that fail the shield's range check, so no writes reach the PLC. This differs from the RQ1 partial-context result (Table~\ref{tab:context_ablation}), where 13~writes \emph{do} execute but 0 succeed: there, the PID correction loop counteracts valid writes faster than the polling window observes. Level Control's resilience thus has two layers: the safety shield blocks out-of-range proposals, and the PID controller corrects in-range writes.

\subsubsection{RQ4: Defense Effectiveness}
\label{sec:eval:rq4}

Table~\ref{tab:defense_comparison} presents the defense comparison across all three scenes using the base Qwen2.5-3B-Instruct model.

% Table 5 --- RQ4 Defense Comparison
% Generated by: python -m cpsforge.analysis.generate_tables
% RQ-specific filtering: limited to 3 runs per (scene, defense) cell
\begin{table}[t]
\centering
\caption{RQ4: Defense comparison for online MITM attacker (base Qwen2.5-3B). Prevention rate = fraction of attack proposals blocked before PLC write. ``None'' = no RQ4 defense; the mandatory safety shield (always active) accounts for any nonzero Blk.\ in ``None'' rows.}
\label{tab:defense_comparison}
\small
\setlength{\tabcolsep}{2pt}
\resizebox{\columnwidth}{!}{%
\begin{tabular}{llrrrrrrrrrr}
\toprule
\textbf{Scene} & \textbf{Defense} & \textbf{Runs} & \textbf{Prop.} & \textbf{Blk.} & \textbf{Wr.} & \textbf{Succ.} & \textbf{ASR\%} & \textbf{Prev.\%} & \textbf{Ph.} & \textbf{Int.} & \textbf{LLM} \\
\midrule
Level Control & None & 3 & 14 & 14 & 0 & 0 & 0.0 & 100.0 & 0 & 0 & 0 \\
 & Baseline & 3 & 0 & 0 & 0 & 0 & 0.0 & 0.0 & 0 & 0 & 0 \\
 & Phase-Aware & 3 & 0 & 0 & 0 & 0 & 0.0 & 0.0 & 0 & 0 & 0 \\
 & Intent & 3 & 18 & 0 & 18 & 3 & 16.7 & 0.0 & 0 & 0 & 0 \\
 & Combined & 3 & 0 & 0 & 0 & 0 & 0.0 & 0.0 & 0 & 0 & 0 \\
 & LLM Defender & 3 & 0 & 0 & 0 & 0 & 0.0 & 0.0 & 0 & 0 & 0 \\
 & LLM+Combined & 3 & 0 & 0 & 0 & 0 & 0.0 & 0.0 & 0 & 0 & 0 \\
\midrule
Sort.~Weight & None & 3 & 13 & 0 & 13 & 7 & 53.8 & 0.0 & 0 & 0 & 0 \\
 & Baseline & 3 & 0 & 0 & 0 & 0 & 0.0 & 0.0 & 0 & 0 & 0 \\
 & Phase-Aware & 3 & 0 & 0 & 0 & 0 & 0.0 & 0.0 & 0 & 0 & 0 \\
 & Intent & 3 & 0 & 0 & 0 & 0 & 0.0 & 0.0 & 0 & 0 & 0 \\
 & Combined & 3 & 0 & 0 & 0 & 0 & 0.0 & 0.0 & 0 & 0 & 0 \\
 & LLM Defender & 3 & 0 & 0 & 0 & 0 & 0.0 & 0.0 & 0 & 0 & 0 \\
 & LLM+Combined & 3 & 0 & 0 & 0 & 0 & 0.0 & 0.0 & 0 & 0 & 0 \\
\midrule
Sort.~Height & None & 3 & 30 & 0 & 30 & 1 & 3.3 & 0.0 & 0 & 0 & 0 \\
 & Baseline & 3 & 30 & 0 & 30 & 4 & 13.3 & 0.0 & 0 & 0 & 0 \\
 & Phase-Aware & 3 & 28 & 0 & 28 & 0 & 0.0 & 0.0 & 0 & 0 & 0 \\
 & Intent & 3 & 27 & 0 & 27 & 2 & 7.4 & 0.0 & 0 & 0 & 0 \\
 & Combined & 3 & 30 & 0 & 30 & 5 & 16.7 & 0.0 & 0 & 0 & 0 \\
 & LLM Defender & 3 & 30 & 30 & 0 & 0 & 0.0 & 100.0 & 0 & 0 & 30 \\
 & LLM+Combined & 3 & 28 & 28 & 0 & 0 & 0.0 & 100.0 & 0 & 0 & 28 \\
\bottomrule
\end{tabular}%
}
\vspace{2pt}
\noindent\footnotesize{3 repeats per configuration; ASR computed post-hoc.}
\end{table}

\textbf{Finding 8: The LLM defender achieves 100\% prevention on Sorting by Height.} The LLM-based anomaly defender blocked all 30 attack proposals (100\% prevention rate) and all 28 proposals in the LLM+Combined configuration. No rule-based defense achieved comparable blocking: baseline, intent, and combined defenses let all writes through to the PLC, where some succeeded (ASR up to 16.7\%). We note that these proposals already passed the safety shield; the LLM defender blocks \emph{shield-legal} writes that rule-based defenses miss. However, a controlled FPR study across all three scenes (Section~\ref{sec:discussion}) reveals that the LLM defender also blocks 95--100\% of legitimate writes (226/229 probes, FPR\,=\,98.7\%), indicating a ``block-everything'' policy rather than discriminating anomaly detection.

\textbf{Finding 9: Phase-aware shield row coincides with zero ASR but does not cause it.} On Sorting by Height, all 28 writes proposed under the phase-aware shield configuration passed through without any blocking (Ph.~=~0, Blk.~=~0), yet none succeeded (ASR~=~0\%). The zero ASR is \emph{not} attributable to the phase-aware shield---it blocks nothing in this configuration. The result most plausibly reflects run-to-run variance: the baseline defense row (which also performs no blocking) achieves 13.3\% ASR across its proposals, so the phase-aware row's 0\% ASR with only 28 proposals lies well within the expected variance range (95\% CI includes 0\% for $n = 28$, $k = 0$). Whether the phase-aware shield influences which proposals the LLM generates---and thereby indirectly affects success---cannot be determined from these results alone.

\textbf{Finding 10: Rule-based defenses do not block LLM-generated writes on Sorting by Height.} On Sorting by Height, baseline, intent, and combined defenses allow all proposed writes through to the PLC without blocking (Blk.=0 in all rows). The Ph.\ and Int.\ columns---which record per-stage block counts---show zero because the phase-aware shield and intent checker evaluate writes but do not block them; their effect manifests through ASR reduction, not write prevention. Only the LLM defender blocks at the proposal level (100\% prevention). Critically, ASR values across defense rows are not directly comparable because the attacker generates different proposals across runs; the relevant comparison is whether any successful attacks occur, not the specific ASR value.

\textbf{Important note on Table~\ref{tab:defense_comparison} ``None'' rows.} The ``None'' defense variant deactivates all RQ4 defense components. The mandatory safety shield is always active; it blocks out-of-range or duration-violating proposals before any RQ4 defense fires. On Level Control ``None'' (RQ4), all 14 proposals are blocked by the safety shield (100\% prevention). This appears to contradict Table~\ref{tab:context_ablation} (RQ1), where Level Control partial context shows 13 writes executed. The difference is explained by run-to-run variation in the LLM's proposed values: the RQ4 ``full context'' attacker proposes different target tags and values than the RQ1 ``partial context'' attacker, and the specific values chosen in the RQ4 runs happened to fall outside the configured safety envelope. Prevention rates in ``None'' rows reflect safety-shield blocking of out-of-range proposals, not RQ4 defense effectiveness.

\textbf{Finding 11: Attacker self-deterrence confounds defense effectiveness measurement on Level Control and Sorting by Weight.} On these two scenes, most defense variants produce 0 attacker proposals---not because the defense blocked them, but because the attacker self-deterred (issued ``wait'' on every LLM call). This limits our ability to attribute 0\% ASR to defense blocking versus attacker passivity. For these scenes, the only interpretable defense comparison rows are those with nonzero proposals (Level Control Intent: 18 proposals, 16.7\% ASR; Sorting by Weight None: 13 proposals, 53.8\% ASR). The defense study is most interpretable on Sorting by Height, where the attacker consistently proposes attacks across all defense configurations. We acknowledge this as a limitation of the full-context evaluation design for RQ4 (see Section~\ref{sec:eval:threats}).

\subsubsection{RQ4 (Finetuned): Defense Against a Fine-Tuned Attacker}
\label{sec:eval:rq4ft}

Table~\ref{tab:defense_finetuned} presents the defense comparison when both attacker and defender use the fine-tuned model (symmetric setup).

% Table RQ4 Finetuned --- Defense Comparison (finetuned attacker)
% Generated by: python -m cpsforge.analysis.generate_tables
% RQ-specific filtering: limited to 3 runs per (scene, defense) cell
\begin{table}[t]
\centering
\caption{RQ4 (finetuned): Defense comparison for finetuned MITM attacker (Qwen2.5-3B+QLoRA, full context). ``None'' = no RQ4 defense; safety shield always active. Sort.~Weight None: 6/13 proposals blocked by safety shield (finetuned model targets \texttt{light\_thresh} with out-of-range values 10.5--15.0). Sort.~Height None: finetuned model self-deterred (0 proposals, full context).}
\label{tab:defense_finetuned}
\small
\setlength{\tabcolsep}{2pt}
\resizebox{\columnwidth}{!}{%
\begin{tabular}{llrrrrrrrrr}
\toprule
\textbf{Scene} & \textbf{Defense} & \textbf{Runs} & \textbf{Prop.} & \textbf{Blk.} & \textbf{Wr.} & \textbf{Succ.} & \textbf{ASR\%} & \textbf{Prev.\%} & \textbf{LLM} \\
\midrule
Level Control & None & 3 & 28 & 0 & 28 & 7 & 25.0 & 0.0 & 0 \\
 & Baseline & 3 & 22 & 0 & 22 & 3 & 13.6 & 0.0 & 0 \\
 & Phase-Aware & 3 & 29 & 0 & 29 & 4 & 13.8 & 0.0 & 0 \\
 & Intent & 3 & 19 & 0 & 19 & 3 & 15.8 & 0.0 & 0 \\
 & Combined & 3 & 21 & 0 & 21 & 3 & 14.3 & 0.0 & 0 \\
 & LLM Defender & 3 & 8 & 8 & 0 & 0 & 0.0 & 100.0 & 8 \\
 & LLM+Combined & 3 & 0 & 0 & 0 & 0 & 0.0 & 0.0 & 0 \\
\midrule
Sort.~Weight & None & 3 & 13 & 6 & 7 & 1 & 14.3 & 46.2 & 0 \\
 & Baseline & 3 & 16 & 16 & 0 & 0 & 0.0 & 100.0 & 0 \\
 & Phase-Aware & 3 & 28 & 18 & 10 & 0 & 0.0 & 64.3 & 0 \\
 & Intent & 3 & 25 & 7 & 18 & 2 & 11.1 & 28.0 & 0 \\
 & Combined & 3 & 30 & 20 & 10 & 1 & 10.0 & 66.7 & 0 \\
 & LLM Defender & 3 & 5 & 5 & 0 & 0 & 0.0 & 100.0 & 0 \\
 & LLM+Combined & 3 & 26 & 26 & 0 & 0 & 0.0 & 100.0 & 0 \\
\midrule
Sort.~Height & None & 3 & 0 & 0 & 0 & 0 & 0.0 & 0.0 & 0 \\
 & Baseline & 3 & 10 & 0 & 10 & 4 & 40.0 & 0.0 & 0 \\
 & Phase-Aware & 3 & 26 & 0 & 26 & 9 & 34.6 & 0.0 & 0 \\
 & Intent & 3 & 26 & 0 & 26 & 8 & 30.8 & 0.0 & 0 \\
 & Combined & 3 & 30 & 0 & 30 & 10 & 33.3 & 0.0 & 0 \\
 & LLM Defender & 3 & 30 & 30 & 0 & 0 & 0.0 & 100.0 & 30 \\
 & LLM+Combined & 3 & 30 & 30 & 0 & 0 & 0.0 & 100.0 & 30 \\
\bottomrule
\end{tabular}%
}
\vspace{2pt}
\noindent\footnotesize{3 repeats per configuration; ASR computed post-hoc.}
\end{table}

\textbf{Finding RQ4-FT-1: LLM defender retains 100\% prevention on Sorting by Height.} On Sorting by Height, the fine-tuned LLM defender blocks all 30 proposals (LLM~=~30), identical to the base model's performance. The fine-tuned attacker achieves 25--40\% ASR against all rule-based defenses on this scene---but the LLM defender stops all of them. For Sorting by Weight, the LLM Defender row shows 100\% prevention (5/5 proposals blocked), but the LLM column is 0---all blocks are attributable to the safety shield, not the LLM defender component: the fine-tuned attacker targets \texttt{light\_thresh} with out-of-range values (10.5--15.0) which the shield rejects before the LLM defender is invoked. Also notable: the Level Control LLM+Combined row shows 0 proposals while LLM Defender alone produces 8; this inversion---fewer proposals under more defenses---is likely run-to-run variance across the 3-repeat design.

\textbf{Finding RQ4-FT-2: Rule-based defenses are weaker against the fine-tuned attacker.} On Sorting by Height, the fine-tuned attacker achieves 40\% ASR against the baseline defense and 30--35\% against phase-aware, intent, and combined---substantially higher than the base model (13.3\%, 0\%, 7.4\%, 16.7\% respectively). The fine-tuned attacker has learned to propose writes that pass rule-based filters.

\textbf{Finding RQ4-FT-3: Fine-tuned attacker self-deterres under no-defense on Sorting by Height.} The fine-tuned model generates zero proposals under no-defense on Sorting by Height, while the base model generates 30 in the same configuration (RQ4 base, Table~\ref{tab:defense_comparison}). This is an unexplained inversion: fine-tuning improves attack capability on some cells (Level Control minimal) while inducing complete self-deterrence on others (Sort.\ Height no-defense). The most plausible explanation is that the fine-tuning data's 83 successful attack examples are biased toward scenes and configurations where defenses were active, so the adapter may have learned to associate attack proposals with defense-present contexts. Testing this hypothesis requires analyzing the fine-tuning examples by configuration, which we leave for future work.

\subsubsection{RQ5: Cross-Model Robustness}
\label{sec:eval:rq5}

Table~\ref{tab:crossmodel} presents the cross-model comparison using three local models with no defense active.

% Table 6 --- RQ5 Cross-Model Comparison
% Generated by: python -m cpsforge.analysis.generate_tables
% RQ-specific filtering: limited to 3 runs per (model, scene) cell
\begin{table}[t]
\centering
\caption{RQ5: Cross-model comparison (no defense). Runs column aggregates across context levels (minimal, full) and attacker types (static LLM, online MITM) for each model--scene pair; ASR is computed over all executed writes within this aggregate. All models run in 4-bit NF4 quantization.}
\label{tab:crossmodel}
\small
\setlength{\tabcolsep}{3pt}
\begin{tabular}{llrrrrr}
\toprule
\textbf{Model} & \textbf{Scene} & \textbf{Runs} & \textbf{Prop.} & \textbf{Writes} & \textbf{Succ.} & \textbf{ASR\%} \\
\midrule
Qwen2.5-3B & Level Control & 12 & 53 & 51 & 7 & 13.7 \\
 & Sort.~Weight & 12 & 0 & 0 & 0 & 0.0 \\
\midrule
Qwen3-1.7B & Level Control & 12 & 58 & 58 & 0 & 0.0 \\
 & Sort.~Weight & 12 & 0 & 0 & 0 & 0.0 \\
\midrule
SmolLM3-3B & Level Control & 12 & 60 & 60 & 7 & 11.7 \\
 & Sort.~Weight & 12 & 0 & 0 & 0 & 0.0 \\
\bottomrule
\end{tabular}
\vspace{2pt}
\noindent\footnotesize{3 repeats per configuration; ASR computed post-hoc.}
\end{table}

\textbf{Finding 12: Attack capability scales with model size.} On Level Control, Qwen2.5-3B achieves 13.7\% ASR, SmolLM3-3B achieves 11.7\%, and Qwen3-1.7B achieves 0\%. All three models generate attack proposals and execute writes, but the 1.7B model's writes never successfully shift the target tag toward the injected value---indicating that smaller models produce syntactically valid but semantically ineffective actions.

\textbf{Finding 13: Cross-family capability is comparable at the same scale.} Qwen2.5-3B (13.7\%) and SmolLM3-3B (11.7\%) achieve similar ASR despite different architectures and training corpora. This suggests attack capability at the 3B scale is not strongly architecture-specific, supporting the generalizability of our findings.

\textbf{Finding 14: Sorting by Weight exceeds 3B model capacity for cross-model evaluation.} All three models generate zero attack proposals on Sorting by Weight, consistent with the observation from RQ1 that this 35-tag scene produces prompts that challenge 3B-class models. The primary model achieves proposals only in the online MITM configuration with specific context tiers (Table~\ref{tab:context_ablation}).

\subsubsection{RQ6: Frontier Model Scaling}
\label{sec:eval:rq6}

RQ6 tests whether a frontier API model (GPT-4o-mini) produces qualitatively different attack outcomes compared to the local 3B-class models evaluated in RQ1--RQ5. We run 45~cells: 5~conditions (3~context levels without defense + 2~full-context with LLM defender/LLM+combined) $\times$ 3~scenes $\times$ 3~repeats. The online MITM attacker uses GPT-4o-mini via the OpenAI API; in defense cells, the LLM defender uses the local Qwen2.5-3B-Instruct, testing cross-model defense generalization.

Table~\ref{tab:frontier} presents the head-to-head comparison. Qwen2.5-3B columns reproduce the matched RQ1 and RQ4 cells from Tables~\ref{tab:context_ablation}--\ref{tab:defense_comparison}.

% Table RQ6 --- Frontier Model Comparison
% Generated by: python scripts/_generate_rq6_table.py
\begin{table*}[t]
\centering
\caption{RQ6: Frontier model comparison---GPT-4o-mini (API) vs.\ Qwen2.5-3B-Instruct (local, 4-bit). No-defense rows (top): online MITM attacker, safety shield always active. Defense rows (bottom): full context, LLM defender uses local Qwen2.5-3B for both models. 3~repeats per cell; ASR\,=\,Succ./Writes computed post-hoc.}
\label{tab:frontier}
\small
\setlength{\tabcolsep}{3pt}
\begin{tabular}{llr|rrrrr|rrrrr}
\toprule
 & & & \multicolumn{5}{c|}{\textbf{GPT-4o-mini}} & \multicolumn{5}{c}{\textbf{Qwen2.5-3B}} \\
\textbf{Scene} & \textbf{Condition} & \textbf{N} & \textbf{Prop.} & \textbf{Wr.} & \textbf{Suc.} & \textbf{ASR\%} & \textbf{Lat.(s)} & \textbf{Prop.} & \textbf{Wr.} & \textbf{Suc.} & \textbf{ASR\%} & \textbf{Lat.(s)} \\
\midrule
Level Control & Minimal & 3 & 30 & 28 & 0 & 0.0 & 2.5 & 23 & 2 & 0 & 0.0 & 13.1 \\
 & Partial & 3 & 30 & 30 & 1 & 3.3 & 3.1 & 20 & 13 & 0 & 0.0 & 18.4 \\
 & Full & 3 & 30 & 30 & 6 & 20.0 & 4.5 & 5 & 0 & 0 & 0.0 & 7.9 \\
\midrule
Sort.~Weight & Minimal & 3 & 30 & 18 & 18 & 100.0 & 2.8 & 13 & 13 & 7 & 53.8 & 17.9 \\
 & Partial & 3 & 30 & 6 & 4 & 66.7 & 2.7 & 12 & 5 & 2 & 40.0 & 19.8 \\
 & Full & 3 & 30 & 0 & 0 & 0.0 & 4.7 & 0 & 0 & 0 & 0.0 & 11.2 \\
\midrule
Sort.~Height & Minimal & 3 & 30 & 30 & 0 & 0.0 & 2.3 & 30 & 30 & 1 & 3.3 & 16.9 \\
 & Partial & 3 & 30 & 30 & 0 & 0.0 & 2.7 & 30 & 25 & 1 & 4.0 & 22.9 \\
 & Full & 3 & 30 & 30 & 0 & 0.0 & 3.8 & 30 & 30 & 4 & 13.3 & 22.3 \\
\midrule
\multicolumn{13}{l}{\textit{Defense active (full context):}} \\
Level Control & LLM Def. & 3 & 30 & 0 & 0 & 0.0 & 5.0 & 0 & 0 & 0 & 0.0 & 24.0 \\
 & LLM+Comb. & 3 & 30 & 0 & 0 & 0.0 & 4.2 & 0 & 0 & 0 & 0.0 & 24.0 \\
Sort.~Weight & LLM Def. & 3 & 30 & 0 & 0 & 0.0 & 5.3 & 0 & 0 & 0 & 0.0 & 24.1 \\
 & LLM+Comb. & 3 & 30 & 0 & 0 & 0.0 & 4.9 & 0 & 0 & 0 & 0.0 & 22.7 \\
Sort.~Height & LLM Def. & 3 & 30 & 0 & 0 & 0.0 & 4.2 & 30 & 0 & 0 & 0.0 & 24.3 \\
 & LLM+Comb. & 3 & 30 & 0 & 0 & 0.0 & 5.0 & 28 & 0 & 0 & 0.0 & 21.6 \\
\bottomrule
\end{tabular}
\vspace{2pt}
\noindent\footnotesize{3 repeats per cell; ASR computed post-hoc. GPT-4o-mini via OpenAI API; Qwen2.5-3B local inference in 4-bit NF4. LLM defender always uses local Qwen2.5-3B-Instruct.}
\end{table*}

\textbf{Finding 15: Self-deterrence is eliminated at frontier-model scale.} GPT-4o-mini generates the maximum 30~proposals (10~per run) in every cell across all 45~runs---zero self-deterrence under any context level or defense configuration. Qwen2.5-3B ranges from 0~proposals (Sort.~Weight full context: complete self-deterrence) to 30 (Sort.~Height: no self-deterrence). This confirms that self-deterrence is a \emph{model-capacity artifact} of 3B-class models, not an inherent property of LLM attackers operating with rich operational context. The implication for defenders is significant: the ``deterrence through transparency'' effect observed in RQ1 cannot be relied upon against larger models.

\textbf{Finding 16: Frontier model achieves the highest ASR in the evaluation on Sorting by Weight.} GPT-4o-mini achieves 100\% ASR (18/18 writes) on Sort.~Weight under minimal context and 66.7\% (4/6) under partial context. These are the highest ASR values in the entire evaluation, exceeding Qwen2.5-3B's 53.8\% and 40.0\% on the same cells. Under full context, GPT-4o-mini generates 30~proposals but all are shield-blocked (0~writes)---the model consistently targets \texttt{heavy\_thresh} with out-of-range values, while under minimal context it targets \texttt{light\_thresh} with shield-legal values. Qwen2.5-3B self-deterres to 0~proposals under full context. The net effect is the same (0\% ASR) but through fundamentally different mechanisms: active blocking vs.\ passive self-suppression.

\textbf{Finding 17: Full context enables frontier-model attacks on Level Control.} GPT-4o-mini achieves 20.0\% ASR (6/30 writes) on Level Control under full context, where Qwen2.5-3B achieves 0\% (5~proposals, all shield-blocked). This is the only configuration in the evaluation where \emph{any} model achieves nonzero ASR on Level Control under full context, and it directly contradicts the self-deterrence hypothesis for frontier-scale models: more context provides \emph{more effective} targeting rather than suppressing attack proposals.

\textbf{Finding 18: Scene-dependent reversal---local model outperforms frontier on Sorting by Height.} On Sort.~Height, GPT-4o-mini achieves 0\% ASR across all context levels despite executing 30/30 writes in every no-defense cell. Qwen2.5-3B achieves 13.3\% ASR under full context (4/30 writes). GPT-4o-mini's writes reach the PLC but produce no observable physical effect---suggesting that the frontier model selects syntactically valid but semantically ineffective targets (e.g., \texttt{enable}, \texttt{conveyor\_entry}) for this discrete sorting scene. The local model's more cautious targeting under full-context guidance yields more effective (though still low) physical impact.

\textbf{Finding 19: LLM defender generalizes across model scale.} In all six defense cells, 0/30 GPT-4o-mini proposals reach the PLC. The local Qwen2.5-3B LLM defender blocks all frontier-model proposals (30/30 on Sort.~Height, 0~writes on all scenes). Compare to Qwen2.5-3B defense cells: on Level Control and Sort.~Weight, the local model self-deterres to 0~proposals so the LLM defender is never tested; on Sort.~Height, 30~proposals are blocked (matching GPT-4o-mini). Cross-model defense generalization holds: a local 3B-class defender blocks a frontier API attacker.

\textbf{Finding 20: Inference latency decreases 4--7$\times$ with API model.} GPT-4o-mini averages 2.3--5.3\,s per LLM call across all cells, compared to 7.9--24.3\,s for local Qwen2.5-3B inference. The latency reduction does not translate to proportionally higher ASR---the bottleneck for attack success is semantic target selection, not inference speed.

\subsubsection{Run-to-Run Variance}

Table~\ref{tab:variance} shows the variance in ASR across repeated runs for key configurations. The LLM defender and phase-aware shield achieve zero variance (ASR~=~0\% across all repeats), demonstrating consistent defense effectiveness. Undefended configurations show moderate variance (std 9--13\,pp), reflecting the stochastic nature of LLM decision-making. The Level Control row aggregates base and finetuned runs (the nonzero mean is entirely attributable to finetuned-model runs; base-model runs achieve 0\% on all Level Control full-context cells). The cross-model comparison confirms that Qwen3-1.7B's zero ASR is not an outlier---all six repeats produce zero successful attacks.

% Table 8 — Variance analysis and confidence intervals
% Generated by: python scripts/compute_variance.py
% Data source: data/raw/live_*/metrics.json, data/raw/campaign_*/metrics.json
%
% STATUS: PARTIAL — 8/17 pairs have ≥2 runs. Run Phase B+C for full CIs.
% Regenerate with: python scripts/compute_variance.py
%
\begin{table}[t]
\centering
\caption{Attack success rate and detector F1 with 95\% confidence intervals from repeated trials. $n$ = number of independent runs. Pairs with $n=1$ show point estimates only.}
\label{tab:variance}
\small
\begin{tabular}{llccc}
\toprule
\textbf{Scene} & \textbf{Attacker} & $n$ & \textbf{ASR (95\% CI)} & \textbf{F1 (95\% CI)} \\
\midrule
Filling Tank & LLM batch & 1 & 0.00 & 0.24 \\
Filling Tank & Random & 1 & 0.20 & 0.81 \\
Filling Tank & Scripted & 2 & 0.17 [0.00, 1.00] & 0.76 [0.76, 0.76] \\
From A to B & LLM batch & 1 & 0.50 & 0.07 \\
From A to B & Random & 3 & 0.23 [0.00, 0.52] & 0.68 [0.68, 0.68] \\
From A to B & Scripted & 5 & 0.27 [0.00, 0.72] & 0.61 [0.24, 0.98] \\
Level Ctrl & Campaign & 1 & 0.22 & 0.00 \\
Level Ctrl & LLM batch & 1 & 0.33 & 0.08 \\
Level Ctrl & Random & 2 & 0.75 [0.11, 1.00] & 0.42 [0.00, 1.00] \\
Level Ctrl & Scripted & 2 & 0.75 [0.00, 1.00] & 0.75 [0.75, 0.75] \\
Sort Height & LLM batch & 1 & 0.00 & 0.05 \\
Sort Height & Random & 2 & 0.06 [0.00, 0.86] & 0.82 [0.81, 0.83] \\
Sort Height & Scripted & 4 & 0.12 [0.00, 0.35] & 0.68 [0.68, 0.68] \\
Sort Weight & Campaign & 2 & 0.39 [0.00, 1.00] & 0.32 [0.00, 1.00] \\
Sort Weight & LLM batch & 1 & 0.67 & 0.21 \\
Sort Weight & Random & 1 & 0.29 & 0.70 \\
Sort Weight & Scripted & 1 & 0.00 & 0.93 \\
\bottomrule
\end{tabular}
\end{table}


\subsubsection{Statistical Power Analysis}
\label{sec:eval:power}

The experiment matrix uses 3~repeats per cell with an attack budget of 10 per run. This design was dictated by resource constraints---each run requires 6--10~minutes of exclusive GPU time with a live PLC---and yields 30--90~total writes per cell depending on attacker behavior. We report Wilson 95\% confidence intervals (CIs) for key ASR results to characterize what the data can and cannot establish.

\textbf{Robust findings} (consistent across all repeats, narrow CIs, or replicated across multiple independent cells):
\begin{itemize}[leftmargin=1.2em, nosep]
    \item \emph{LLM defender 100\% blocking on Sort.~Height}: 0~successes in 30~writes under base model, 0~in 30~under finetuned model, 0~in 30~under GPT-4o-mini (three independent cells). Wilson 95\% CI for ASR: $[0\%, 11.6\%]$ per cell. Zero variance across all repeats (Table~\ref{tab:variance}).
    \item \emph{GPT-4o-mini self-deterrence elimination}: 30/30~proposals in all 45~runs, zero variance. Structural, not stochastic.
    \item \emph{GPT-4o-mini 100\% ASR on Sort.~Weight minimal}: 18/18~writes. Wilson 95\% CI: $[82\%, 100\%]$. However, this is computed over writes that passed the safety shield (18~of 30~proposals); among all proposals, 12/30 were shield-blocked.
    \item \emph{Qwen3-1.7B 0\% ASR}: 0~successes in 58~writes across 12~runs. Wilson 95\% CI: $[0\%, 6.2\%]$. Zero variance across 6~per-run estimates.
    \item \emph{Full-context self-deterrence on Sort.~Weight}: 0~proposals generated across 3~repeats per cell for Qwen2.5-3B, replicated in both base and finetuned models (2~independent cells). GPT-4o-mini does \emph{not} self-deterre (30~proposals, all shield-blocked), confirming model-capacity dependence.
    \item \emph{Static LLM zero proposals}: 0~proposals across all 9~scene$\times$repeat cells (27~runs). Structural, not stochastic.
\end{itemize}

\textbf{Moderate-confidence findings} (consistent direction, but CIs overlap zero or are wide):
\begin{itemize}[leftmargin=1.2em, nosep]
    \item \emph{GPT-4o-mini 20\% ASR on Level Control full}: 6/30 writes. Wilson 95\% CI: $[9\%, 38\%]$. Excludes zero; a real effect, but magnitude uncertain.
    \item \emph{GPT-4o-mini 66.7\% ASR on Sort.~Weight partial}: 4/6 writes. Wilson 95\% CI: $[30\%, 90\%]$. Very wide due to small $n$; directionally above Qwen (40\%, 2/5 writes) but not separable.
    \item \emph{Sort.~Weight MITM 53.8\% ASR}: 7/13 writes. Wilson 95\% CI: $[28\%, 78\%]$. The point estimate is above chance but the CI is wide.
    \item \emph{Fine-tuning improves Level Control minimal}: 11/30 = 36.7\% (FT) vs.\ 0/2 = 0\% (base). Wilson CI for FT: $[22\%, 55\%]$. The base cell has only 2~executed writes, limiting comparison power.
    \item \emph{3B $>$ 1.7B on Level Control}: Qwen2.5-3B 7/51 = 13.7\%, CI $[7\%, 26\%]$; SmolLM3-3B 7/60 = 11.7\%, CI $[6\%, 22\%]$; Qwen3-1.7B 0/58, CI $[0\%, 6\%]$. The two 3B models are indistinguishable from each other; both are consistently above 1.7B.
\end{itemize}

\textbf{Preliminary findings} (single-cell observations or high variance):
\begin{itemize}[leftmargin=1.2em, nosep]
    \item \emph{Sort.~Height full context 13.3\% ASR}: 4/30 writes. Wilson 95\% CI: $[5\%, 30\%]$. Run-to-run std = 13.1\,pp (Table~\ref{tab:variance}), with individual runs ranging from 0\% to 40\%.
    \item \emph{Random outperforms MITM on Sort.~Height}: 4/30 = 13.3\% vs.\ 1/30 = 3.3\%. Wilson CIs overlap: $[5\%, 30\%]$ vs.\ $[1\%, 17\%]$. Not statistically separable at $n=30$.
    \item \emph{FT attacker 40\% ASR on Sort.~Height Baseline defense}: 4/10 writes. Wilson CI: $[17\%, 69\%]$. Very wide; this cell should be treated as directional evidence only.
\end{itemize}

With 3~repeats per cell, the design has adequate power to detect large effects ($\Delta$ASR $> 30$\,pp) and zero-vs-nonzero differences, but cannot reliably distinguish moderate effects ($\Delta$ASR $\approx 10$\,pp) from noise. We report all findings with this caveat and recommend that future work use $\geq$10~repeats per cell to narrow CIs below $\pm$15\,pp.

\subsection{Threats to Validity}

\textbf{Internal validity.} Attack success is measured post-hoc by comparing tag values before and after each executed write. This approximation may undercount attacks where the PLC controller rapidly corrects the deviation within one polling cycle ($<$500\,ms). The 3-repeat design provides limited statistical power for rare-event metrics. For a cell with $k$ successes out of $n$ writes, the Wilson 95\% confidence interval on ASR can exceed $\pm$30\,pp for $n \leq 10$. All quantitative findings should be interpreted as point estimates with this uncertainty in mind; cells with $n \geq 10$ and $k \geq 1$ provide more reliable estimates (e.g., Sort.~Weight None: 7/13 = 53.8\%, 95\% CI $[28\%, 77\%]$). The variance table (Table~\ref{tab:variance}) provides run-to-run dispersion data for key configurations.
\label{sec:eval:threats}

\textbf{External validity.} All experiments use one PLC model (Siemens S7-1200) with software-emulated plants (Factory~I/O). Results may not transfer to other PLC families or physical plants. The 3B-class models studied represent the lower end of current LLM capability; larger models may produce qualitatively different attack strategies.

\textbf{Construct validity.} The ``attack success'' metric measures tag displacement toward the injected value, not physical process impact (e.g., tank overflow, misrouted parts). A successful tag shift may not always cause operational harm, and conversely, subtle shifts below the detection threshold may accumulate over time.
